% !TEX root = ../main.tex
% Agent documentation:
% Purpose: identify conversion barriers across the customer journey for a new
% drone-delivery category.
% Downstream use: Section 07a turns these barriers into tactics for CAC,
% conversion rate, and average order value.
% Revision guardrails: keep each journey step tied to a behavioral mechanism or
% documented risk. When goals in Section 05 change, check the final summary for
% outdated goal references.
\subsection{Kundereise}

Dronelevering er en ukjent kategori for norske forbrukere. Uten
erfaring å lene seg på blir hvert steg i kjøpsprosessen mer bevisst enn for
en Foodora-bestilling, og opplevd risiko bremser der rutine ellers ville
drevet beslutningen. Zipline må derfor aktivt styre valgarkitekturen gjennom
hele kundereisen.

Analysen følger Kotler og Kellers femstegs kjøpsprosess
\parencite{kotlerkeller2016} og bruker valgarkitekturteknikker fra
\textcite{muenscher2016} der de har konkrete implikasjoner.

\subsubsection{Steg 1: Problemerkjennelse}

Foodora og Wolt har normalisert 30 minutters leveringstid, og høy adopsjon
viser at forbrukerne aksepterer dette som godt nok. Gapet mellom ønsket og
faktisk tilstand er for lite til å
utløse problemerkjennelse av seg selv. Zipline må derfor reframe
referansepunktet fra leveringstid til kvalitet ved ankomst, altså fra ``30
minutter er raskt nok'' til ``maten bør ankomme varm''. Reframingen treffer
hardest i tidskritiske situasjoner der ventetid har konkrete konsekvenser,
som glemte ingredienser under matlaging eller gjester på vei. Kontekstuell
annonsering rettet mot slike situasjoner aktiverer problemerkjennelsen der
tidsfordelen merkes tydeligst.

\subsubsection{Steg 2: Informasjonssøking}

Når kategorien er ny, lener forbrukerne seg på personlige kilder fremfor
massemedia \parencite{rogers2003}. Synlig bruk blant naboer og anbefalinger
fra teknologi- eller matprofiler har dermed større effekt enn betalt
annonsering. Droner har samtidig en større iboende synlighet enn budbaserte
tjenester, og sprer kjennskap som biprodukt av selve leveringen.
Synligheten har en bakside, fordi støy og visuell tilstedeværelse er blant
de vanligste innvendingene mot droner. Operasjonene bør derfor likevel være så
diskrete som mulig.

\subsubsection{Steg 3: Evaluering av alternativer}

Opplevd risiko er hovedbarrieren. Dronelevering reiser spørsmål om
sikkerhet, personvern og pålitelighet som ikke finnes for budbasert
levering, og noen forbrukere vil avvise tjenesten uavhengig av pris og
hastighet. Zipline kan møte dette på to måter. Den første er å forenkle
sikkerhetsinformasjonen ved å oversette tekniske spesifikasjoner til
gjenkjennelige tall. ``2 millioner leveranser uten alvorlige hendelser''
 overbeviser mer enn informasjon om redundanssystemer. 
Den andre er en
gratis eller sterkt rabattert førstegangsleveranse, fordi opplevd risiko
reduseres mest effektivt gjennom egen erfaring.

Fraværet av byttekostnader (jf.\ kraft 3 i Porters femkraftsanalyse) gjør at
``Foodora eller Zipline'' stilles på nytt ved hver bestilling. Pris- og
hastighetsfortrinnene fra posisjoneringsanalysen må derfor kommuniseres i
selve bestillingsøyeblikket, ikke bare i kjennskapsfasen.

\subsubsection{Steg 4: Kjøpsbeslutningen}

Valgarkitektur har størst effekt i bestillingsøyeblikket, og Ziplines app bør
utformes for å styre valgprosessen. Tre teknikker fra
\textcite{muenscher2016} er direkte anvendbare. \textit{Standardvalg}
innebærer at dronelevering er forhåndsvalgt i områder med dekning, fordi
forbrukere i stor grad aksepterer forhåndsinnstillinger. \textit{Redusert
friksjon} gjennom én-klikks gjenbestilling fjerner søke- og evalueringssteget
for gjentakende brukere. \textit{Ankring} oppnås ved å vise leveringsgebyret
ved siden av konkurrentenes typiske gebyr, slik at prisfordelen blir konkret
i beslutningsøyeblikket. Tiltakene forutsetter bestilling gjennom Ziplines
egen plattform. Dersom levering også tilbys gjennom tredjepartsapper, er
kontrollen over valgarkitekturen begrenset.

\subsubsection{Steg 5: Etterkjøpsvurdering og gjenkjøp}

Etterkjøpsvurderingen avgjør om engangskunden blir gjentakskunde.
\textcite{kahnemanriis2005} viser at hukommelsen domineres av det mest
intense øyeblikket (peak) og avslutningen (end). For Zipline er
leveringsøyeblikket det naturlige peak-punktet og bør designes for å
forsterke opplevelsen av presisjon, for eksempel gjennom sanntidssporing av
de siste meterne og en notifikasjon når pakken er klar for henting.
Produktkvalitet ved ankomst er end og siste inntrykk.

Uten byttekostnader (jf.\ steg 3) må gjenkjøp aktivt bygges som vane.
Påminnelser ved typiske bestillingstidspunkt holder Zipline fremme i
bevisstheten og konkurrerer med impulsen til å åpne Foodora. Et abonnement
med fast rabatt skaper byttekostnader i en bransje som ellers mangler dem,
og flytter konkurransen fra enkeltkjøp til perioder.

\subsubsection{Oppsummering: konverteringsutfordringer langs kundereisen}

Foodora og Wolt taper kunder primært mellom vurdering og kjøp (pris,
utvalg, leveringstid). Zipline risikerer å tape dem tidligere, mellom
oppmerksomhet og vurdering, fordi potensielle kunder kjenner til tjenesten
men aldri vurderer den når de mangler erfaring og oppfatter risikoen som
høy.
De første 90 dagene etter lansering handler derfor primært om å senke
opplevd risiko og få folk til å prøve tjenesten gjennom
sikkerhetskommunikasjon og lav terskel for førstegangsbruk. Når de første
kundene er konvertert, kan synlig bruk og anbefalinger fra eksisterende
kunder hjelpe med tillitsbygging hos nye potensielle kunder.
